\documentclass[11pt]{article}
\usepackage[margin=1in]{geometry}
\usepackage{amsmath,amssymb,amsthm,mathtools}
\usepackage{graphicx}
\usepackage{microtype}
\usepackage[hidelinks]{hyperref}
\usepackage{xcolor}

\newtheorem{theorem}{Theorem}
\newtheorem{lemma}{Lemma}
\newcommand{\R}{\mathbb{R}}
\newcommand{\G}{\mathbb{G}}
\newcommand{\tr}{\operatorname{tr}}
\newcommand{\sgn}{\operatorname{sgn}}
\newcommand{\Span}{\operatorname{span}}

\title{A Two-Plane Counterexample to the Atomic Condition\\
for a Pl\"ucker-Coordinate $\ell^p$ Integrand}
\author{Automated mathematics run \texttt{fa\_banach\_001}, lane 6}
\date{19 July 2026}

\begin{document}
\maketitle

\noindent\textbf{Status.} Candidate counterexample, likely valid; prepared for
expert review.  The result disproves the all-exponent atomic-condition claim
suggested in Remark~6.4 of De Rosa--Tione~\cite{DRT}.

\section{Source question}

The source is Antonio De Rosa and Riccardo Tione, \emph{Regularity for graphs
with bounded anisotropic mean curvature}, arXiv:2011.09922v2, published in
\emph{Inventiones Mathematicae} 230 (2022), 463--507.  On page 26,
Remark~6.4 reports numerical evidence that the integrand on $\G(4,2)$ induced
by the $\ell^p$ norm of the six Pl\"ucker coordinates satisfies (SAC), hence
(AC), for every $p\in(1,\infty)$, and states that no analytic proof was found.

\begin{center}
\includegraphics[width=\textwidth]{figures/open_problem_crop.png}
\end{center}

The precise sentence is: ``numerical simulations indicated that for
$\mathcal G(\cdot)=\|\cdot\|_{\ell^p}$, $p\in(1,\infty)$,
$\Psi_{\mathcal G}$ fulfills (SAC), hence actually satisfies (AC)''~\cite[Remark
6.4]{DRT}.

\section{Definitions}

Write the coordinates of a bivector $z\in\Lambda^2(\R^4)$ as
\[
 z=(z_{12},z_{13},z_{14},z_{23},z_{24},z_{34}).
\]
It is simple precisely when
\[
 z_{12}z_{34}-z_{13}z_{24}+z_{14}z_{23}=0.
\]
Every unit simple bivector determines an unoriented plane in $\G(4,2)$,
with the two choices $z$ and $-z$ giving the same plane.

For $p>1$, let
\[
 \mathcal G_p(z)=\left(\sum_{i<j}|z_{ij}|^p\right)^{1/p},
 \qquad
 \Psi_p(T_z)=\mathcal G_p(z)
\]
for unit simple $z$.  If $K(z)$ is the skew-symmetric matrix with
$K(z)_{ij}=z_{ij}$ for $i<j$, and
$\phi_p(s)=\sgn(s)|s|^{p-1}$, equation (6.2) of~\cite{DRT} gives
\begin{equation}\label{eq:B}
 B_p(T_z)=
 \frac{K(\phi_p(z))K(z)^{\mathsf T}}{\mathcal G_p(z)^{p-1}}.
\end{equation}
For a probability measure $\mu$ on $\G(4,2)$, put
\[
 A_p(\mu)=\int_{\G(4,2)}B_p(T)\,d\mu(T).
\]
Here the normal dimension is $n=2$.  Condition (AC2) says that if
$\dim\ker A_p(\mu)=2$, then $\mu$ must be a Dirac mass~\cite[Definition
3.1]{DRT}.

\section{Proof intuition}

The source proves that every averaged matrix has rank at least two, so a
failure of (AC) must occur exactly at rank two with a nontrivial measure.  We
choose two simple bivectors obtained from one another by interchanging a large
and a small Pl\"ucker coordinate.  Their $B_p$ matrices have complementary
block forms.  The equal average is symmetric and has only two eigenvalues,
each repeated twice.  As $p$ varies from the nonsmooth endpoint $1$ to the
Euclidean exponent $2$, the smaller eigenvalue changes sign.  At an interior
exponent it is therefore zero with multiplicity two, giving exactly the
forbidden non-Dirac equality case.

\section{Counterexample theorem}

\begin{theorem}\label{thm:main}
There exists $p_0\in(1,2)$ such that the Pl\"ucker-coordinate $\ell^{p_0}$
integrand $\Psi_{p_0}$ on $\G(4,2)$ does not satisfy the atomic condition.
More precisely, there are two distinct planes $T_x,T_y\in\G(4,2)$ such that,
for
\[
 \mu=\tfrac12(\delta_{T_x}+\delta_{T_y}),
\]
one has $\dim\ker A_{p_0}(\mu)=2$.
\end{theorem}

\begin{proof}
Fix $a,b>0$ such that
\[
 a+b=1,\qquad a\ne b,
\]
and set $t=\sqrt{ab/2}$.  In Pl\"ucker-coordinate order define
\begin{equation}\label{eq:xy}
 x=(a,-t,t,t,t,-b),\qquad
 y=(b,-t,t,t,t,-a).
\end{equation}
Both are simple, because
\[
 x_{12}x_{34}-x_{13}x_{24}+x_{14}x_{23}
 =-ab+t^2+t^2=0,
\]
and the same calculation applies to $y$.  They are unit bivectors since
\[
 |x|^2=|y|^2=a^2+b^2+4t^2=a^2+b^2+2ab=(a+b)^2=1.
\]
The planes $T_x$ and $T_y$ are distinct: if they agreed, their unit simple
representatives would satisfy $y=\pm x$, which is impossible when $a\ne b$
and $t>0$.

Fix $p>1$ temporarily and abbreviate
\[
 \alpha=a^{p-1},\qquad \beta=b^{p-1},\qquad
 \theta=t^{p-1},\qquad
 \rho_p=\mathcal G_p(x)^{p-1}=\mathcal G_p(y)^{p-1}>0.
\]
The equality defining $\rho_p$ follows because $x$ and $y$ have the same
multiset of absolute coordinates.  Applying~\eqref{eq:B} gives
\begin{equation}\label{eq:Bx}
 \rho_p B_p(T_x)=
 \begin{pmatrix}
 d_1&0&-r&-r\\
 0&d_1&-r&r\\
 -s&-s&d_2&0\\
 -s&s&0&d_2
 \end{pmatrix},
\end{equation}
and
\begin{equation}\label{eq:By}
 \rho_p B_p(T_y)=
 \begin{pmatrix}
 d_2&0&-s&-s\\
 0&d_2&-s&s\\
 -r&-r&d_1&0\\
 -r&r&0&d_1
 \end{pmatrix},
\end{equation}
where
\[
 d_1=\alpha a+2\theta t,\quad
 d_2=\beta b+2\theta t,\quad
 r=\alpha t+\theta b,\quad
 s=\beta t+\theta a.
\]
Consequently
\begin{equation}\label{eq:average}
 A_p(\mu)=\frac1{\rho_p}
 \begin{pmatrix}
 d&0&-c&-c\\
 0&d&-c&c\\
 -c&-c&d&0\\
 -c&c&0&d
 \end{pmatrix},
\end{equation}
with
\begin{equation}\label{eq:dc}
 d=\frac{\alpha a+\beta b}{2}+2\theta t,
 \qquad
 c=\frac{(\alpha+\beta)t+\theta}{2};
\end{equation}
we used $a+b=1$ in the formula for $c$.

The characteristic polynomial of the unnormalized matrix in
\eqref{eq:average} is
\[
 \big((\lambda-d)^2-2c^2\big)^2.
\]
Thus its eigenvalues are
\[
 d-\sqrt2c\quad\text{and}\quad d+\sqrt2c,
\]
each with multiplicity two.  Define $D(p)=d-\sqrt2c$ for $p>1$ and extend it
continuously to $p=1$.  Since $\alpha,\beta,\theta\to1$ as $p\downarrow1$,
\begin{equation}\label{eq:D1}
 D(1)=\frac12+2t-\sqrt2\left(t+\frac12\right).
\end{equation}
The right side is strictly increasing as a function of $t$, and it vanishes
at $t=1/(2\sqrt2)$.  But $a\ne b$ and $a+b=1$ imply
$ab<1/4$, hence $t<1/(2\sqrt2)$.  Therefore $D(1)<0$.

At $p=2$, one has $\alpha=a$, $\beta=b$, and $\theta=t$, so
\[
 d=\frac{a^2+b^2}{2}+2t^2=\frac12,
 \qquad c=t,
\]
and hence
\begin{equation}\label{eq:D2}
 D(2)=\frac12-\sqrt2t=\frac12-\sqrt{ab}>0.
\end{equation}
By the intermediate value theorem there exists $p_0\in(1,2)$ such that
$D(p_0)=0$.  At this exponent $d=\sqrt2c>0$, so the matrix in
\eqref{eq:average} has eigenvalues $0,0,2d,2d$.  The positive normalization
$\rho_{p_0}$ does not change the kernel.  Thus
\[
 \dim\ker A_{p_0}(\mu)=2,
\]
although $\mu$ is supported on the two distinct planes $T_x,T_y$ and is not
Dirac.  This violates (AC2), proving the theorem.
\end{proof}

\paragraph{Concrete numerical location.}
Taking $a=4/5$ and $b=1/5$, bisection of the explicit continuous function
$D$ gives
\[
 p_0\approx1.1593840621.
\]
This decimal is a reproducibility check only; the proof uses the strict signs
in~\eqref{eq:D1}--\eqref{eq:D2}, not numerical computation.

\section{Verification and limitations}

The script \texttt{code/verify\_counterexample.py} independently reconstructs
the matrices from~\eqref{eq:B}, checks the Pl\"ucker and norm identities,
factors the characteristic polynomial, and bisects the concrete root.  The
computational output is not used as a proof.

The theorem gives a complete negative answer to the universal all-$p$ claim
suggested in Remark~6.4.  It does not classify the exponents for which (AC)
holds.  It also does not contradict Theorem~6.1 of~\cite{DRT}: that theorem
establishes (AC1) for every $p\in(1,\infty)$, whereas the construction violates
only the equality-rigidity clause (AC2).

\section{Bounded novelty check and review recommendation}

On 19 July 2026, the run's four lightweight indexes had no match for
arXiv:2011.09922 or the core scalar-atomic/Pl\"ucker $\ell^p$ terms.  Primary
arXiv searches used the exact phrases ``scalar atomic condition'', ``uniform
scalar atomic condition'', and ``Pl\"ucker atomic condition'', together with
the source title and authors.  No prior resolution of Remark~6.4 was found.
The closest later primary sources inspected were arXiv:2401.10858v2 (a
necessary strict-polyconvexity condition for (AC)) and arXiv:2601.10647v2
(codimension-one $\ell^p$ anisotropies in a different Michael--Simon result);
neither treats this two-plane $\G(4,2)$ obstruction.  This bounded search is
not an exhaustive novelty guarantee.

Human review should first compare~\eqref{eq:B} with equation (6.2) of the
source, then audit the signs in~\eqref{eq:Bx}--\eqref{eq:By}.  The remaining
kernel and intermediate-value argument is elementary.

\begin{thebibliography}{9}
\bibitem{DRT}
A. De Rosa and R. Tione,
\emph{Regularity for graphs with bounded anisotropic mean curvature},
Invent. Math. 230 (2022), 463--507;
arXiv:2011.09922v2.
\end{thebibliography}

\end{document}
